\begin{question}

Two observables $A1$ and $A2$, which do not involve time explicitly,
are known not to commute,

\begin{equation*}
  [A1 , A2] \ne 0,
\end{equation*}

yet we also know that $A1$ and $A2$ both commute with the Hamiltonian:

\begin{equation*}
  [A1, H] = 0, [A2, H] = 0.
\end{equation*}

\begin{questionpart}
  Prove that the energy eigenstates are, in general, degenerate.
\end{questionpart}

\begin{questionpart}
  Are there exceptions?
\end{questionpart}

\begin{questionpart}
  As an example, you may think of the central-force problem $H =
  \vectorbold{p}^2/2m + V(r)$, with $A1 \rightarrow L_z$, $A2
  \rightarrow L_x$.
\end{questionpart}

\end{question}

\begin{purpose}
  OK, this problem is meant to be our first foray into wave mechanics
  and there are a LOT of things going on here.  This is where we start
  going from abstract operators to real-world integrals. I fell flat
  on my face the first time I tried to attempt this one until some
  external help allowed me to perform a rectal craniotomy and muddle
  through it.  Let's start by listing a few of the things going on
  here:

  \begin{itemize}

  \item Contrary to what you may recall from equation 1.33 in the
    \textit{Griffiths} book, $\hat{p} \ne \left( -i \hbar
    \frac{\partial}{\partial x}\right)$.  Rather, working from the
    foundation of Hamltonian mechanics found in Goldstein, Sakurai
    posits the existence of a momentum operator $\hat{p}$ and then
    shows how it must behave.  In doing so, he develops a generalized
    expression of that behavior in equation 1.249.  If you try to
    treat the $-i \hbar \ldots$ as the operator itself, then you'll
    get to equation 1.248 and start cursing at the world saying ``How
    in the actual fsck did you move that $\frac{\partial}{\partial x}$
    into the middle of the dyad while multiplying by
    1?!\@\#\@!?!!1!eleven!1?  We'll cover this one soon.

  \item There is a difference between an abstract operator, say
    $\hat{p}$, and the expression of that operator in an integral.
    Considering the example of $\hat{p}$ when acting on a ket holding
    the position, the proper expression in physical space is $-i \hbar
    \left(\frac{\partial}{\partial x}\right)$ (well, more properly
    it's $-i \hbar \nabla \vectorbold{x}$) but if you're acting on a
    ket that represents the momentum, then it gets much more
    complicated and becomes $p$, if you're operating on a ket that
    represents Schr\"{o}dinger's cat (the sadist...who kills
    kittens!?) then I'm pretty sure it spontaneously generates a black
    hole.

  \item It's probably worth reviewing Goldstein chapter XXX where he
    derives the generators from Hamiltonian mechanics.
  \end{itemize}
\end{purpose}

\begin{answer}
  
\end{answer}
